\problema{Cheapest Flights Within K Stops}{787}{src/01-grafos/787-cheapest-flights-k-stops.cpp}{M}

Hay \texttt{n} ciudades numeradas de \texttt{0} a \texttt{n-1} y una lista
\texttt{flights} de vuelos \texttt{[from, to, price]} (aristas dirigidas y
ponderadas, sin pesos negativos). Dados \texttt{src}, \texttt{dst} y
\texttt{k}, buscamos el precio más barato de \texttt{src} a \texttt{dst}
usando \textbf{a lo más \texttt{k} escalas} (paradas intermedias), o
\texttt{-1} si no existe tal ruta. Usar a lo más \texttt{k} escalas equivale
a usar a lo más \texttt{k+1} vuelos (aristas) en el camino.

\paragraph{La idea, con un ejemplo de la vida real.} Quieres volar de una
ciudad a otra lo más barato posible, pero con un límite: como mucho
\texttt{k} \emph{escalas} (paradas intermedias). Una ruta con muchas escalas
puede salir más barata, pero si se pasa del límite no te sirve, aunque cueste
menos.

Podríamos pensar en Dijkstra (el del problema anterior), pero Dijkstra solo
se fija en el precio total y no lleva la cuenta de \emph{cuántos vuelos}
usaste. Y aquí eso importa muchísimo: dos rutas que cuestan igual no valen lo
mismo si una usa más escalas que la otra. Necesitamos un método que avance
\emph{un vuelo a la vez} y nos deje detenernos justo después de \texttt{k+1}
vuelos.

\begin{center}
\ttfamily
\begin{tabular}{c}
n=4, flights=[[0,1,100],[1,2,100],[2,0,100],[1,3,600],[2,3,200]]\\[6pt]
0 --100--> 1 --100--> 2 --200--> 3\\
\phantom{0 --100--> }1 --------600--------> 3
\end{tabular}
\end{center}

Con \texttt{src=0, dst=3, k=1}: la ruta $0\to1\to2\to3$ cuesta
$100+100+200=400$ pero usa \emph{2} escalas (nodos 1 y 2), más de las
permitidas. La única ruta con a lo más 1 escala es $0\to1\to3$, que cuesta
$100+600=700$. Respuesta: \texttt{700}.

\paragraph{Cómo lo resuelve: Bellman-Ford por rondas (patrón nuevo).}
\emph{Bellman-Ford} es otro algoritmo de caminos. Su idea es sencilla: hace
\emph{rondas}, y en cada ronda revisa \textbf{todos} los vuelos e intenta
abaratar el precio de llegar a cada ciudad. Lo bonito es esto: después de $r$
rondas, el precio apuntado para cada ciudad es lo más barato que se puede
llegar \emph{usando a lo más $r$ vuelos}. Como $k$ escalas son $k+1$ vuelos,
hacemos \textbf{exactamente $k+1$ rondas}: ni una más (usaríamos más escalas
de las permitidas) ni una menos (nos quedaríamos cortos).

\paragraph{El truco de la copia (la parte clave).} En cada ronda trabajamos
sobre una \emph{fotografía congelada} de los precios de la ronda anterior. Es
decir: construimos un arreglo nuevo, \texttt{nuevaDist}, mirando siempre los
precios ``de ayer'' (\texttt{dist}), y solo al terminar toda la ronda
reemplazamos los viejos por los nuevos.

¿Por qué? Si actualizáramos los precios ``en el momento'' (como hace
Dijkstra), una sola ronda podría encadenar varios vuelos seguidos: abaratar
$0\to1$ y, en esa misma ronda, usar ya el nuevo precio de $1$ para abaratar
$1\to2$. Eso metería dos vuelos (dos escalas) de contrabando en una sola
ronda, y daría una ruta más barata pero \emph{tramposa}, porque usa más
escalas de las permitidas. Al copiar la fotografía, cada ronda solo puede
sumar \emph{exactamente un} vuelo por camino, sin importar en qué orden
revisemos la lista de vuelos.

\lstinputlisting{src/01-grafos/787-cheapest-flights-k-stops.cpp}

\complejidad{$O(k\cdot E)$}{$O(n)$}

\paragraph{¿Qué significa esa complejidad?} $O(k\cdot E)$ significa que
hacemos $k+1$ rondas y en cada ronda miramos las $E$ conexiones (los vuelos)
una vez; es decir, el trabajo es simplemente ``número de rondas por número de
vuelos''. El espacio $O(n)$ es porque solo guardamos un precio por cada una
de las $n$ ciudades (más su copia).

\paragraph{Consejos para la entrevista (Amazon).} Este problema es la puerta
de entrada al patrón de \textbf{Bellman-Ford con rondas limitadas}, muy útil
cuando el límite no es sobre el costo sino sobre el \emph{número de aristas}
del camino (aquí, las escalas). El error más común es olvidar la copia y
actualizar sobre el mismo arreglo, lo que deja colar rutas con más escalas de
las permitidas; Amazon suele preguntar justo por qué esa copia es necesaria.
Otro tropiezo es confundir ``k escalas'' con ``k vuelos'': hay que hacer
\texttt{k+1} rondas, no \texttt{k}. Casos raros: \texttt{src == dst} (cuesta
\texttt{0}, ni siquiera vuelas), destino imposible de alcanzar con cualquier
número de escalas (\texttt{-1}) y una \texttt{k} tan grande que el límite ya
no estorba (el algoritmo simplemente termina de calcular el más barato de
todos, sin que las rondas de más dañen el resultado).
